\subsection{Actual closed projective subvarieties in the normalized model}
\label{subsec:native-closed-projective-spaces}

\begin{theorem}[Genuine closed-projective normalized comparison]
\label{projective:native-closed-normalized-comparison}
\lean{EnvelopeProperty.ProjBerk.normalizedClosedLocus,
  EnvelopeProperty.ProjBerk.isClosed_normalizedClosedLocus,
  EnvelopeProperty.ProjBerk.nativeClosedProjectiveHomeomorph,
  EnvelopeProperty.ProjBerk.nativeClosedProjectiveHomeomorph_apply,
  EnvelopeProperty.ProjBerk.nativeClosedProjectiveHomeomorph_support_toIdeal}
\leanok
\uses{scheme:relative-native-closed-immersion,projective:normalized-native-homeomorphism}
Let $P$ be the polynomial projective scheme
\[
 P=\operatorname{Proj}(R[X_0,\ldots,X_n]),
\]
and let $g:X\to P$ be an actual closed immersion of schemes. For any prescribed real seminorm $v$ on $R$, the
native valued-point space of $X$ over the true composite base morphism
is homeomorphic to the normalized fixed-base projective points whose
actual homogeneous prime support lies in the genuine scheme image of
$g$. This normalized locus is closed. The entire native image point and
its full original polynomial seminorm kernel are preserved.
\end{theorem}
\begin{proof}
The actual scheme image is closed. The normalized support map is genuinely
continuous, so its inverse image is a true closed normalized locus.
The proved prescribed-base native closed-immersion homeomorphism
identifies the native points of $X$ with the actual support-closed locus
in native projective space. The genuine global normalized/native
homeomorphism preserves precisely that prime support, so restricting it
to this true locus gives another genuine homeomorphism. Compose the
first map with the inverse of this restriction.

The support-image homeomorphism retains the original whole native map,
and the second restricted homeomorphism has genuine inverse identities.
Their composition therefore preserves the entire native image point,
not only its support or coordinate norms. The already proved genuine
native prime-support identity identifies that point's entire homogeneous
prime ideal with the full polynomial valuation kernel of the resulting
normalized point. No presentation, chart or point-comparison certificate
is assumed.
\end{proof}

\begin{theorem}[Actual closed-projective compactness and closed-subscheme specialization]
\label{projective:native-closed-projective-compactness}
\lean{EnvelopeProperty.ProjBerk.compactSpace_nativeClosedProjective,
  EnvelopeProperty.ProjBerk.t2Space_nativeClosedProjective,
  EnvelopeProperty.ProjBerk.isCompact_nativeClosedProjective_relativeLocus,
  EnvelopeProperty.ProjBerk.nativeProjectiveSubschemeHomeomorph,
  EnvelopeProperty.ProjBerk.nativeProjectiveSubschemeHomeomorph_apply,
  EnvelopeProperty.ProjBerk.compactSpace_nativeProjectiveSubscheme}
\leanok
\uses{projective:native-closed-normalized-comparison,projective:native-fixed-base-compactness,scheme:relative-native-closed-immersion}
The actual prescribed-base native space of $X$ is compact and Hausdorff,
and its whole fixed-base locus is a compact subset of the absolute native
space. These assertions hold for arbitrary $R$ and $v$, without any
uniform coefficient bound. For every actual ideal-sheaf datum $I$ on
polynomial projective space, its true closed subscheme has the resulting
homeomorphism to normalized points supported on the genuine support of
$I$, retaining the entire original inclusion map, and is compact.
\end{theorem}
\begin{proof}
Actual native projective space above $v$ is genuinely compact without
bounding all coefficient values. The proved relative native closed
embedding of $g$ has closed image there, and its genuine subspace
identification transports compactness to its source. Equivalently the
preceding homeomorphism identifies the source with a closed subset of
the compact normalized fixed-base space. That normalized space is
Hausdorff, so the genuine embedding also gives Hausdorffness of the
source. The compact-subtype criterion yields compactness of its original
fixed-base locus in the absolute native point space.

For ideal-sheaf data, Mathlib constructs the actual closed subscheme and
its actual closed immersion, with scheme image precisely the datum's
true support. Apply the preceding comparison and compactness producers
to that inclusion. The proved image/support identity rephrases the
normalized locus as the stated support locus; the full native inclusion
point remains unchanged.

The proofs retain nilpotents, empty closed subschemes and $n=0$, and need
no field, completeness, reducedness, finite-generation or noetherian
assumption. No nonemptiness of a closed subscheme is asserted. This is a
point-space and topology comparison, not an analytic structure-sheaf or
coherent analytic ideal identification.
\end{proof}
